% MI20091C
% Thoraxschmerz, beschwerdefrei
% 14.0
% 2013-11-30

\label{m:akutes-koronarsyndrom-beschwerdefrei}
%
\begin{itemize}
    \item Lagerung: Oberkörper hoch
  \item Bewegungensverbot, Schonung
  \item Beengende Kleidung öffnen
    \item \ce{O2}-Gabe gemäß \FullPrettyRef{m:sauerstoffberieselung}
  \item Bestmögliches Monitoring
  \item \II{Nitro-Spray}: Manche Patienten haben einen
  Nitroglyzerin-Spray (z.\,B. Nitrolingual{\texttrademark})
  zur Einnahme bei Symptomen eines Akuten Koronarsyndroms
  verschrieben bekommen. Der Spray soll weiter \E{vom
  Patienten selbstständig} in der \E{verschriebenen
  Dosierung} genommen werden, solange der systolische
  Blutdurck RR\textsubscript{syst}\,{\textgreater}\,100\MmHg*
  ist. Der Helfer muss sich versichern, dass das Medikament
  für den jeweiligen Patienten verschrieben wurde und dass
  die vorgegebene Dosierung eingehalten wird. Eine
  regelmäßige Blutdruckkontrolle ist notwendig.
  \item Transportentscheidung: Abt. f. Innere Medizin
\end{itemize}